\documentclass[11pt,a4paper]{article}
\usepackage[utf8]{inputenc}
\usepackage[T1]{fontenc}
\usepackage[ngerman,english]{babel}
\usepackage{lmodern}
\usepackage{microtype}
\usepackage{geometry}
\geometry{margin=2.6cm}
\usepackage{amsmath,amssymb,amsthm,mathtools}
\usepackage{booktabs,longtable,enumitem}
\usepackage{hyperref}
\usepackage[nameinlink,noabbrev]{cleveref}
\usepackage[numbers,sort&compress]{natbib}

\hypersetup{colorlinks=true,linkcolor=black,citecolor=black,urlcolor=black}
\newtheorem{theorem}{Theorem}[section]
\newtheorem{reduction}[theorem]{Reduction Theorem}
\theoremstyle{definition}
\newtheorem{definition}[theorem]{Definition}

\newcommand{\Hur}{\mathcal H}
\newcommand{\Sp}{\mathcal S_p}
\newcommand{\Dp}{\mathcal D_p}
\newcommand{\Op}{\mathcal O_p}
\newcommand{\Def}{\mathscr D}
\newcommand{\Z}{\mathbb Z}
\newcommand{\Q}{\mathbb Q}

\title{\textbf{Local Obstructions and Global Diophantine Geometry\\
in the Defect-Protection Property of the Bamberg Model}\\[0.4em]
\large Lokale Obstruktionen und globale diophantische Geometrie\\
bei der Defektschutz-Eigenschaft des Bamberger Modells}
\author{Thomas Hoffbauer\\Bamberg, Germany}
\date{\today}

\begin{document}
\maketitle

\begin{abstract}
We study a low-gap defect-protection property in an arithmetic model based on Hurwitz quaternions for primes \(p\equiv7\pmod8\). On each Hurwitz norm shell we distinguish a diagonal sector \(D\) and a non-diagonal sector \(O\), with minimal positive defects \(g_D(p)\) and \(g_O(p)\). We prove
\[
g_D(p)=4\Longrightarrow g_O(p)\ge4,\qquad
g_D(p)=9\Longrightarrow g_O(p)\ge9,
\]
using divisibility, archimedean product--sum estimates, parity constraints and local obstructions modulo \(3,7,8,16\). The next admissible layer \(g_D=25\) is not protected: at \(p=31\), \(g_D(31)=25\) and \(g_O(31)=20\). For \(g_D=36\), the problem is reduced to six explicit global Diophantine cases. The first residual case is further reduced to a genus-one quartic, coupled Pell equations and a square problem in two Lucas-type sequences.
\end{abstract}

\selectlanguage{ngerman}
\section*{Deutsche Fassung}

\section{Einleitung und Problemstellung}
Wir untersuchen eine Defektschutz-Eigenschaft in einem auf Hurwitz-Quaternionen basierenden arithmetischen Modell für Primzahlen
\[
p\equiv7\pmod8.
\]
Auf jeder Hurwitz-Normschale unterscheiden wir einen diagonalen Sektor \(D\) und einen nichtdiagonalen Sektor \(O\). Die minimalen positiven Defekte werden mit \(g_D(p)\) bzw. \(g_O(p)\) bezeichnet. Für \(k>0\) definieren wir
\[
L_k:\qquad g_D(p)=k\Longrightarrow g_O(p)\ge k.
\]

\section{Hurwitz-Schalen und Defektfunktion}
Sei \(\Hur\) die Hurwitz-Ordnung in der rationalen Hamiltonschen Quaternionenalgebra. Für eine ungerade Primzahl \(p\) betrachten wir
\[
\Sp=\{q\in\Hur:N(q)=p\}.
\]
Wir verwenden verdoppelte Koordinaten
\[
q=\frac12(A+Bi+Cj+Dk),
\]
wobei \(A,B,C,D\in\Z\) alle dieselbe Parität besitzen und
\[
A^2+B^2+C^2+D^2=4p.
\]

\begin{definition}[Defekt]
\[
\Def(A,B,C,D)=\frac{A}{2}\left[B(A^2-2p)+ACD\right].
\]
\end{definition}

\begin{definition}[Sektoren]
\[
\Dp=\{(A,B,C,D)\in\Sp:|A|=|B|\},
\quad
\Op=\{(A,B,C,D)\in\Sp:|A|\ne|B|\}.
\]
Ferner
\[
g_D(p)=\min_{\Dp,\Def\ne0}|\Def|,
\qquad
g_O(p)=\min_{\Op,\Def\ne0}|\Def|.
\]
\end{definition}

\section{Diagonale Reduktion}
Mit
\[
s=\frac{C+D}{2},\qquad r=\frac{C-D}{2}
\]
gilt im halbzahlig-diagonalen Ast
\[
2p=A^2+s^2+r^2,
\qquad
|\Def|=A^2r^2.
\]
Für \(p\equiv7\pmod8\) ergibt die mod-\(16\)-Analyse unter \(g_D<100\)
\[
\boxed{g_D\in\{4,9,25,36\}.}
\]
Der Wert \(25\) ist nicht geschützt:
\[
g_D(31)=25,\qquad g_O(31)=20.
\]

\section{Schutzsatz \(L_4\)}
Für \(g_D=4\) gilt
\[
2p=s^2+5.
\]
Im ganzzahligen Hurwitz-Zweig ist \(4\mid\Def\). Im halbzahligen Zweig gilt \(A\mid\Def\). Die Fälle \(|\Def|=1,2,3\) werden durch Teilbarkeit, Produkt--Summen-Abschätzungen und lokale Kongruenzen ausgeschlossen. Der letzte Ast reduziert sich auf
\[
z^2-3k^2=1,\qquad v^2-2z^2=-3,
\]
was modulo \(3\) unmöglich ist.

\begin{theorem}
Für jede Primzahl \(p\equiv7\pmod8\) gilt
\[
\boxed{g_D(p)=4\Longrightarrow g_O(p)\ge4.}
\]
\end{theorem}

\section{Schutzsatz \(L_9\)}
Für \(g_D=9\) gilt
\[
2p=s^2+10,\qquad s\equiv2\pmod4.
\]
Der ganzzahlige Zweig erlaubt unterhalb \(9\) nur \(|\Def|=4,8\). Im halbzahligen Zweig bleiben nach Teilbarkeit und einer uniformen Abschätzung nur
\[
(|\Def|,|A|)=(3,3),(5,5),(7,7).
\]
Diese drei Fälle werden durch Schranken für \(|B|\), die Substitution
\[
v=\frac{C+D}{2},\qquad z=\frac{C-D}{2},
\]
und lokale Obstruktionen modulo \(3,7,8\) eliminiert. Die ganzzahligen Fälle \(|\Def|=4,8\) werden separat ausgeschlossen.

\begin{theorem}
Für jede Primzahl \(p\equiv7\pmod8\) gilt
\[
\boxed{g_D(p)=9\Longrightarrow g_O(p)\ge9.}
\]
\end{theorem}

\section{Die Schicht \(g_D=36\)}
Es treten zwei Diagonalfamilien auf:
\[
\mathrm{D36}_{13}:\quad 2p=s^2+13,
\qquad
\mathrm{D36}_{37}:\quad 2p=s^2+37.
\]
Der Minimalitätsfilter gegen \(g_D\in\{4,9,25\}\) folgt aus
\[
(t-s)(t+s)=\lambda-\mu.
\]
Damit erhält man
\[
s^2\ge81\quad(\lambda=13),
\qquad
s^2\ge121\quad(\lambda=37).
\]

Sei
\[
m=|\Def|<36,\quad a=|A|,\quad b=|B|,\quad m=aq.
\]
Die notwendige archimedische Ungleichung lautet
\[
(b-a)x+b(\lambda-a^2)-a\lambda+\frac{a^3}{2}
+\frac{ab^2}{2}-2q\le0,
\qquad x=s^2.
\]
Unter den scharfen Untergrenzen für \(x\) überlebt kein Fall mit \(b>a\). Da \(b=a\) diagonal wäre,
\[
\boxed{1\le b<a,\qquad a\ge3.}
\]
Mit \(v=(C+D)/2\) und \(z=(C-D)/2\) folgt die universelle binäre Gleichung
\[
\boxed{
(a-b)v^2-(a+b)z^2
=
\frac{b(b^2-a^2)}2\pm2q.
}
\]

\section{Exakte Reduktion von \(L_{36}\)}
Teiler-, Paritäts- und archimedische Filter liefern \(372\) strukturierte Tupel \((\lambda,m,a,b)\). Wegen der beiden Vorzeichen entstehen \(744\) Vorzeichenäste. Lokale Filter reduzieren
\[
744\longrightarrow252\longrightarrow21\longrightarrow6.
\]

\begin{longtable}{@{}lll@{}}
\toprule
\((\lambda,m,a,b)\) & Binäre Gleichung & Quadratbedingung\\
\midrule
\endhead
\((13,15,5,1)\) & \(2v^2-3z^2=-3\) & \(s^2=v^2+z^2\)\\
\((13,25,25,1)\) & \(12v^2-13z^2=-157\) & \(s^2=v^2+z^2+300\)\\
\((13,33,3,1)\) & \(v^2-2z^2=9\) & \(s^2=v^2+z^2-8\)\\
\((37,15,3,1)\) & \(v^2-2z^2=-7\) & \(s^2=v^2+z^2-32\)\\
\((37,25,25,23)\) & \(v^2-24z^2=-551\) & \(s^2=v^2+z^2+540\)\\
\((37,33,3,1)\) & \(v^2-2z^2=9\) & \(s^2=v^2+z^2-32\)\\
\bottomrule
\end{longtable}

\begin{reduction}
Die Aussage
\[
g_D(p)=36\Longrightarrow g_O(p)\ge36
\]
ist im vorliegenden Modell auf sechs explizite globale diophantische Restfälle reduziert.
\end{reduction}

\section{Globaler Restfall 1}
Aus
\[
2v^2-3z^2=-3,\qquad s^2=v^2+z^2
\]
folgt mit \(v=3u\)
\[
z^2-6u^2=1,\qquad s^2-15u^2=1.
\]
Eine rationale Parametrisierung der ersten Konik liefert
\[
y^2=36t^4+48t^2+1.
\]
Die projektive Glättung besitzt Geschlecht \(1\).

Schreibe
\[
t=\frac rq,\qquad \gcd(r,q)=1,\qquad d=q^2-6r^2.
\]
Integrabilität erzwingt \(d\mid12\). Lokale Filter eliminieren \(d=6,-2,3\). Die beiden verbleibenden Normklassen \(d=1,-6\) führen auf dieselbe Quartik
\[
\boxed{\mathcal C:\quad y^2=360x^4+60x^2+1}
\]
mit Pell-Nebenbedingung
\[
\boxed{X^2-6x^2=1.}
\]

\section{Lucas-Reduktion}
Aus dem gekoppelten System folgt
\[
5W^2-2y^2=3,\qquad W=12x^2+1.
\]
Nach Zerlegung der positiven Lösungen in zwei Einheitenorbits und Restriktion auf \(W\equiv1\pmod{12}\) erhält man
\[
A_0=1,\quad A_1=1177,
\qquad
B_0=265,\quad B_1=382129,
\]
mit
\[
\boxed{T_{k+2}=1442T_{k+1}-T_k.}
\]
Offen bleibt, ob
\[
\frac{A_k-1}{12}
\]
für \(k\ge1\) oder
\[
\frac{B_k-1}{12}
\]
für \(k\ge0\) ein Quadrat sein kann.

\section{Forschungsstand}
\[
\boxed{\mathrm{L4:\ PROVED}},
\qquad
\boxed{\mathrm{L9:\ PROVED}},
\]
\[
\boxed{\mathrm{L36:\ REDUCED\ TO\ 6\ GLOBAL\ DIOPHANTINE\ CASES}.}
\]

\section{Offene Probleme}
\begin{enumerate}
\item Bestimmung aller Pell-kompatiblen Integralpunkte von
\[
y^2=360x^4+60x^2+1.
\]
\item Entscheidung der beiden Quadratprobleme in den Folgen \(A_k,B_k\).
\item Globale Analyse der fünf weiteren L36-Restfälle.
\item Erst danach vollständige Entscheidung von \(L_{36}\).
\end{enumerate}

\clearpage
\selectlanguage{english}
\section*{English Version}

\section{Introduction}
We study a defect-protection property in an arithmetic model based on Hurwitz quaternions for primes \(p\equiv7\pmod8\). On each Hurwitz norm shell we distinguish a diagonal sector \(D\) and a non-diagonal sector \(O\), with minimal positive defects \(g_D(p)\) and \(g_O(p)\). Define
\[
L_k:\qquad g_D(p)=k\Longrightarrow g_O(p)\ge k.
\]

\section{Hurwitz shells and the defect}
Let \(\Hur\) denote the Hurwitz order. For an odd prime \(p\),
\[
\Sp=\{q\in\Hur:N(q)=p\}.
\]
Write
\[
q=\frac12(A+Bi+Cj+Dk),
\]
where \(A,B,C,D\in\Z\) have common parity and
\[
A^2+B^2+C^2+D^2=4p.
\]
Define
\[
\Def(A,B,C,D)=\frac A2\left[B(A^2-2p)+ACD\right].
\]
The diagonal and non-diagonal sectors are given by \(|A|=|B|\) and \(|A|\ne|B|\), respectively.

\section{Diagonal reduction}
With
\[
s=\frac{C+D}{2},\qquad r=\frac{C-D}{2},
\]
the half-integral diagonal branch satisfies
\[
2p=A^2+s^2+r^2,\qquad |\Def|=A^2r^2.
\]
For \(p\equiv7\pmod8\) and \(g_D<100\),
\[
\boxed{g_D\in\{4,9,25,36\}.}
\]
At \(p=31\),
\[
g_D(31)=25,\qquad g_O(31)=20.
\]

\section{Protection at \(g_D=4\)}
For \(g_D=4\),
\[
2p=s^2+5.
\]
The integral branch satisfies \(4\mid\Def\). In the half-integral branch, \(A\mid\Def\). The cases \(|\Def|=1,2,3\) are excluded by divisibility, product--sum inequalities and local congruences. The last branch reduces to
\[
z^2-3k^2=1,\qquad v^2-2z^2=-3,
\]
which is impossible modulo \(3\).

\begin{theorem}
For every prime \(p\equiv7\pmod8\),
\[
\boxed{g_D(p)=4\Longrightarrow g_O(p)\ge4.}
\]
\end{theorem}

\section{Protection at \(g_D=9\)}
For \(g_D=9\),
\[
2p=s^2+10,\qquad s\equiv2\pmod4.
\]
The half-integral branch reduces to the three cases
\[
(|\Def|,|A|)=(3,3),(5,5),(7,7),
\]
which are eliminated by bounds on \(|B|\), the substitution \(v=(C+D)/2,\ z=(C-D)/2\), and local obstructions modulo \(3,7,8\). The integral cases \(|\Def|=4,8\) are excluded separately.

\begin{theorem}
For every prime \(p\equiv7\pmod8\),
\[
\boxed{g_D(p)=9\Longrightarrow g_O(p)\ge9.}
\]
\end{theorem}

\section{The \(g_D=36\) layer}
Two diagonal families occur:
\[
\mathrm{D36}_{13}:\quad 2p=s^2+13,
\qquad
\mathrm{D36}_{37}:\quad 2p=s^2+37.
\]
A minimality filter against \(g_D\in\{4,9,25\}\) yields
\[
s^2\ge81\quad(\lambda=13),
\qquad
s^2\ge121\quad(\lambda=37).
\]
For \(m=|\Def|<36\), \(a=|A|\), \(b=|B|\), \(m=aq\), the necessary inequality
\[
(b-a)x+b(\lambda-a^2)-a\lambda+\frac{a^3}{2}
+\frac{ab^2}{2}-2q\le0
\]
forces
\[
1\le b<a,\qquad a\ge3.
\]
Every surviving branch satisfies
\[
\boxed{
(a-b)v^2-(a+b)z^2
=
\frac{b(b^2-a^2)}2\pm2q.
}
\]

\section{Exact reduction of \(L_{36}\)}
The divisibility, parity and archimedean filters leave \(372\) structured tuples. The two signs give \(744\) signed branches. Local sieving produces
\[
744\longrightarrow252\longrightarrow21\longrightarrow6.
\]
The six remaining cases are the same six cases listed in the German section above.

\begin{reduction}
Within the present model,
\[
g_D(p)=36\Longrightarrow g_O(p)\ge36
\]
is reduced to six explicit global Diophantine cases.
\end{reduction}

\section{Global residual case 1}
The first case reduces to
\[
z^2-6u^2=1,\qquad s^2-15u^2=1,
\]
then to the genus-one quartic
\[
y^2=36t^4+48t^2+1.
\]
After denominator reduction, the two remaining norm classes lead to the common quartic
\[
\boxed{\mathcal C:\quad y^2=360x^4+60x^2+1}
\]
with Pell side condition
\[
\boxed{X^2-6x^2=1.}
\]

\section{Lucas reduction}
Elimination gives
\[
5W^2-2y^2=3,\qquad W=12x^2+1.
\]
The relevant terms split into two sequences
\[
A_0=1,\ A_1=1177,
\qquad
B_0=265,\ B_1=382129,
\]
both satisfying
\[
\boxed{T_{k+2}=1442T_{k+1}-T_k.}
\]
The remaining question is whether \((A_k-1)/12\) or \((B_k-1)/12\) can be a square in the stated index ranges.

\section{Current status}
\[
\boxed{\mathrm{L4:\ PROVED}},
\qquad
\boxed{\mathrm{L9:\ PROVED}},
\]
\[
\boxed{\mathrm{L36:\ REDUCED\ TO\ 6\ GLOBAL\ DIOPHANTINE\ CASES}.}
\]

\section{Open problems}
\begin{enumerate}
\item Determine all Pell-compatible integral points on
\[
y^2=360x^4+60x^2+1.
\]
\item Decide the two associated Lucas square problems.
\item Solve the remaining five global \(L_{36}\) cases.
\item Only then decide \(L_{36}\) itself.
\end{enumerate}

\bibliographystyle{plainnat}
\bibliography{references}
\end{document}
